\documentclass[11pt]{article}
\usepackage[a4paper,margin=2.35cm]{geometry}
\usepackage{fontspec}
\usepackage{amsmath,amssymb}
\setmainfont[ItalicFont=Noto Serif Italic]{Malgun Gothic}
\setsansfont{Malgun Gothic}
\setlength{\parindent}{1.2em}
\setlength{\parskip}{0.45em}
\setlength{\emergencystretch}{3em}
\begin{document}

$K$를 $P(\varepsilon)$에 포함된 차수 $2^n$의 순환체라 하자. $K$ 역시 분해체임을 주장한다.\footnote{원분체의 부분체로서 차수 $2^n$의 순환 분해체를 택한 것은 하세의 예를 본뜬 것이다.} $K$를 기초체로 할 때 사원수체의 지수를 $m$이라 하면, $m=1$ 또는 $m=2$이다. 그런데 $K$에 대한 상대차수 $(p-1)/2^n$가 홀수인 분해체 $P(\varepsilon)$가 존재하므로, 두 번째 경우는 불가능하다. 따라서 $m=1$이다. 이는 곧 $K$가 분해체라는 뜻이다.

따라서 무한히 많은 $n$에 대하여, 극소 분해체인 차수 $2^n$의 순환체가 존재한다.

\end{document}
